\chapter{Registro de pruebas y evidencias}
\pendiente{Usar una ficha por prueba con identificador, objetivo, entorno, pasos, resultado esperado, resultado observado, criterio de éxito, incidencias y archivos de evidencia. Referenciar cada ficha desde el texto principal.}
\section{Matriz de cumplimiento}
\begingroup\small
\noindent\begin{tabularx}{\textwidth}{>{\raggedright\arraybackslash}p{0.47\textwidth}>{\raggedright\arraybackslash}X}
\toprule
Requisito & Evidencias pendientes\\
\midrule
Varias aplicaciones y dos soluciones utilizadas & Instalación, cifrado, descifrado y comparación\\
Clave PGP publicada & Huella y recuperación desde servidor público\\
Intercambio con compañeros & Envío y recepción con descifrado comprobado\\
Linux instalado con LUKS & Particionado, instalación y verificación de arranque\\
Memoria en Poliformat & Confirmación de subida sin cifrar\\
Memoria enviada por correo cifrada & Destinatarios autorizados y prueba de envío\\
\bottomrule
\end{tabularx}
\endgroup
\chapter{Material complementario}
\pendiente{Añadir comandos comentados, diagramas, capturas pertinentes y glosario cuando aporten valor. Mantener las explicaciones y resultados esenciales en los capítulos principales. No adjuntar claves privadas, frases de paso ni material de recuperación secreto.}
